\section{A subquadratic width polynomial under a two-sided density bound}
\label{sec:localized-results}
The preceding theorem needs only an upper density bound. If the input
law has density bounded away from zero as well, localization gives a
stronger guarantee for the same shallow ReLU learner. Throughout this
section assume
\[
 0<p_-\le p_X(x)\le p_+<\infty\quad(-1\le x\le1).
\]
Initially $|Y|\le1$. Write $f_0(x)=\E[Y\mid X=x]$,
$R_0=\E(Y-f_0(X))^2$, and $A(g)=R(g)-R_0$.
The comparison class $\mathcal G_C$ contains all shallow width-$C$
ReLU functions with arbitrary finite coefficients. All conclusions
below control the posterior-average raw risk and therefore the
ordinary raw posterior mean. The prior and observations are unchanged.
Only the likelihood scale differs from the evidence calibration used
in the preceding sections.

\begin{theorem}[A localized width-only Gaussian oracle]
\label{thm:localized-width}
Let $1\le C\le n$ be a declared comparison-width budget. There is a
deterministic likelihood scale depending on $n,C,\delta$ and the stated
density bounds, but not on the comparison coefficients or its risk,
such that, with probability at least $1-\delta$, simultaneously for
every $N\ge6C+6$ and $0<\nu\le1$,
\begin{equation}
 R(\bar z_S)\le\Pi_SR(z)
 \le\inf_{g\in\mathcal G_C}R(g)
       +\widetilde O\left((C+1)^{2/5}n^{-3/10}\right).
 \label{eq:localized-budget}
\end{equation}
In the nontrivial regime, put
\[
 e_{n,C}=\widetilde c(C+1)^{2/5}n^{-3/10},\qquad
 \lambda\asymp\left(\frac{e_{n,C}}{C+1}\right)^2,
\]
with sufficiently large fixed logarithmic factors. The physical
Gaussian observation variance is $n\lambda\nu/N$.
A sufficient sample size for excess risk $\varepsilon$ is
\begin{equation}
 n\ge\widetilde O\left((C+1)^{4/3}\varepsilon^{-10/3}\right).
 \label{eq:localized-samples}
\end{equation}
Thus the width exponent at fixed accuracy is $4/3$. Any prescribed
$C_n=n^u$, $u<3/4$, gives a vanishing additive remainder with the
single deterministic scale
$\lambda_n=\widetilde\Theta(n^{-3/5-6u/5})$.
\end{theorem}

\begin{corollary}[Schedules that do not use the comparison width]
\label{cor:localized-universal}
One sufficiently large fixed power $k$ in
$\lambda=n^{-3/2}\log^k(n/\delta)$ gives, on a common event for all
comparison widths satisfying $(C+1)\sqrt\lambda\le1$, all
$N\ge6C+6$, and all $0<\nu\le1$,
\begin{align}
 R(\bar z_S)\le\Pi_SR(z)
 \le{}&R_0+\{1+O(1/\log n)\}\inf_{g\in\mathcal G_C}A(g)\\
 &+\widetilde O((C+1)n^{-3/4}+n^{-1/3})+O(1/\log n).
 \label{eq:localized-universal}
\end{align}
Hence every fixed growth $C_n=n^u$, $u<3/4$, is allowed by one
schedule without knowing $C$. Its additive logarithmic term is
weaker in accuracy than the declared-budget polynomial bound.
A different universal schedule,
$\lambda=\widetilde\Theta(n^{-3/5})$, gives the coefficient-one bound
\[
 R(\bar z_S)\le\Pi_SR(z)\le\inf_{g\in\mathcal G_C}R(g)
       +\widetilde O\left((C+1)n^{-3/10}
                          +(C+1)^2n^{-9/10}\right).
\]
\end{corollary}

The bounds are asymptotic statements in the displayed nontrivial
regimes; the proof gives explicit sufficiently large logarithmic
powers and sample conditions. A fixed $\lambda=1$ handles the finitely
many excluded small sample sizes by the same Gaussian energy bound.
There is no mixture over comparison widths or likelihood scales.
The declared-budget schedule in Theorem~\ref{thm:localized-width} uses
$C$; the schedules in Corollary~\ref{cor:localized-universal} do not.

\paragraph{Proof structure.}
The spline compiler in Lemma~\ref{lem:spline-path} has bounded output,
risk increment $O(Ct)$ and path variation $O((C+1)/t)$. The new step
controls empirical quadratic norms through the total variation of the
second derivative, rather than a global output envelope.
Section~\ref{sec:sharp-design} proves this norm comparison with
remainder $\widetilde O(\eta^{-5}n^{-3}P^2)$.
Section~\ref{sec:localized-proof} combines it with localized multiplier
bounds and the exact Gaussian posterior energy identity. Substitution
of the stated schedules proves the theorem and corollary. The common
sample events are function-class events independent of $N,\nu$.
For each declared $C$, choose a deterministic approximate risk minimizer;
a summable confidence allocation over the integer budgets makes the
claims simultaneous, changing only logarithms.

\begin{corollary}[Marginally sub-Gaussian labels]
\label{cor:localized-subgaussian}
Suppose instead that $\E\exp(Y^2/s_Y^2)\le2$ for a fixed known upper
bound $s_Y$. The coefficient-one conclusions
\eqref{eq:localized-budget}--\eqref{eq:localized-samples} and the second
conclusion of Corollary~\ref{cor:localized-universal} hold with additional
logarithmic factors and constants depending on $s_Y$. The first
conclusion holds in the possibly weaker form
\[
 R(\bar z_S)\le\Pi_SR(z)
 \le\{1+O(1/\log n)\}\inf_{g\in\mathcal G_C}R(g)
       +\widetilde O((C+1)n^{-3/4}+n^{-1/3})+O(1/\log n).
\]
Observed labels and the returned predictions remain unrestricted.
The physical likelihood scale is $\lambda=M_n^2\lambda'$, where
$M_n=\max\{1,s_Y\sqrt{c\log(n/\delta)}\}$ and $\lambda'$ is the
corresponding bounded-label schedule, with enlarged logarithmic factors.
\end{corollary}
\begin{proof}
Set $Y_M=\clip(Y,[-M_n,M_n])$ only in the analysis. Its joint law
with $X$ has exactly the original input density. With probability
at least $1-\delta/4$, every training label already lies in this
interval, so the actual posterior is exactly the posterior for the
clipped-label sample. Analyze its observable $z/M_n$ and labels
$Y_M/M_n$. The parameter prior stays unchanged. Its path variation is
at most $Q/2$, its prior empirical risk is at most two, and its Gaussian
energy identity is the same as in Section~\ref{sec:localized-proof}.
A normalized bounded spline comparator $q$ is implemented physically
as $M_nq$; its balanced energy is $2M_nP(q)$. The extra factor $M_n$
and the Gaussian comparison precision cost are logarithmic in $n/\delta$.
Thus every step of the localized argument has the same powers of $n,C$,
with enlarged logarithmic factors. Multiplication by $M_n^2$ converts
normalized risks back to physical risks without changing the leading
coefficient. For the universal schedule, enlarge the fixed logarithmic
power further so that its normalized design remainder is at most
$1/(M_n^2\log n)$; its physical contribution is then still
$O(1/\log n)$. The multiplicative approximation slack stays
$O(1/\log n)$ under risk rescaling.

For completeness, this is a transfer of population risks as well as
training posteriors. The deterministic bound
$\Pi_SQ^2\le C(1+(\lambda')^{-2})$ holds on the clipped training sample.
Since $|z(x)|\le Q/\sqrt2$,
\[
 |\Pi_SR_Y(z)-\Pi_SR_{Y_M}(z)|
 \le C(1+(\lambda')^{-1})\{\E(Y-Y_M)^2\}^{1/2}
           +\E|Y^2-Y_M^2|.
\]
The right side is negligible for a sufficiently large fixed $c$ in
$M_n$, uniformly over the learner width and scale. No law is conditioned
on a label event. For an arbitrary raw comparison function $g$, the
pointwise projection inequality
\[
 (\clip(g,[-M_n,M_n])-Y)^2
 \le(g-Y)^2+(|Y|-M_n)_+^2
\]
first bounds its clipped risk by $R(g)$ plus a label-tail term.
Changing that bounded comparator's labels from $Y$ to $Y_M$ costs
at most $2M_n\E|Y-Y_M|+\E|Y^2-Y_M^2|$.
The clipped comparator still has at most $3C+2$ affine breakpoints.
These observations justify the same coefficient-free compiler and
risk infimum. A union with the clipped-label statistical event and
the all-training-labels event completes the proof.
\end{proof}
